\part{Unresolved comparisons and current gates}
\section{Statements not asserted as proved}

There is no unchecked theorem silently used in the preceding part.
The implications with hypotheses, including B8 and C15's universality
reduction, are proved implications; their hypotheses are not supplied
by this run. The following are the remaining identifications.

\paragraph{G1: the exact iterated character-independent functional.}
The common oscillator kernel, A1's normalized zero orbit, B11's
restriction seesaw, A5's finite-frame identity and A7's actual single
theta lift have been established. The recovered \texttt{MainX} is
a stronger uniform relation-preservation statement for an iterated
family. To apply it, one must construct its actual common geometry,
tuples, archimedean maps, definite class kernels and arithmetic links,
and identify its normalized functional on the selected 032 vectors.
The concrete all-place tower fails the canonical node-gap and
signature requirements. The abstract two-step case is more flexible:
with global Witt index zero, one can take $m=2$ and the Bessel
unipotent rank zero. Its middle space must be totally definite.
A7's natural multiplicity space is totally definite precisely when
the selected first-list CM types agree; that direct realization then
gives a divisor-generated balanced line. A different source
realization remains a live possibility. The last two-lift kernel
and covolume identity, rather than an intermediate hyperbolic tower,
is the current matching problem.

\paragraph{G2: turning the explicit certificates into a uniform
arithmetic lower bound for the prescribed iterated family.}
Q1--Q10 and B5--B7 give actual formulas under specified inputs.
The finite unipotent search yields $|J|\geq1/(4S_{\rm ext})$ with
an explicit grid size after its Gaussian envelope and lattice basis
are supplied. An arithmetic Gram bound needs a coefficient field,
denominator and bounds at all embeddings. Q9 provides resultant
denominators for specified integral Hecke matrices. The remaining
step is the normalization and overlap data of the particular
iterated family in G1. Positivity alone cannot supply it, as B4
and B6 show. No universal cycle-degree estimate is asserted.

\paragraph{G3: an exceptional non-CM Hodge summand.}
C16 supplies infinitely many non-CM source factors, and C13 gives
an explicit algebraic projector onto a Weil space. One nonzero
projected cycle would generate that two-dimensional rational space.
The known curve-product maps in C12 and C16 are annihilated by this
projector when the target has the generic Weil-type divisor algebra,
by C14. The remaining task is a specified map or correspondence
with nonzero Weil projection, or the full-rank pairing data in C04.
C15's characteristic-zero universality implication is a source
reduction, not a proof of that conjecture or of the extraction gate.
The countability result restricts fixed arithmetic towers over
$\C$; it does not preclude the dense algebraic families constructed
here or a further deformation argument.

\paragraph{G4: the chosen integral correspondence and connection.}
C17 identifies the exact horizontal correspondence and local-model
data needed to compare the 032 image cycle with a Kudla--Rapoport
special cycle. C06--C08 prevent replacing those data by an
undifferentiated positivity claim. C11 proves the sharp integral
height-two derivative estimate. C18 supplies actual integral tame
theta quotients and isolates the common sigma connection for
nonzero-degree tensors. One must now match the archive's harmonic
covariant derivatives, period normalization and line-bundle degree
to those functions, and identify the horizontal correspondence on
the desired integral model. Neither a raw analytic frame nor a
vertical local intersection is silently used in that role.

These gates are recorded with dated changes in \texttt{LEDGER.md}.
The mandated three-hour progress review has not been replaced by
this mathematical inventory. Its timing is controlled by the actual
run clock and recorded separately.
